\section{Analisis Keamanan}

\subsection{Langkah Mitigasi}

\paragraph{Kebocoran kunci privat.}
Berkas kunci privat dibuat dengan izin \texttt{0600} pada berkas sementara di direktori
yang sama, lalu di-\textit{rename}, sehingga isi kunci tidak pernah terbaca pihak lain dan
berkas lama tidak rusak bila penulisan gagal. Berkas tidak ditimpa tanpa opsi
\texttt{force} yang eksplisit. Modul kunci tidak mencatat isi kunci pada log. Pada
container, \texttt{.dockerignore} menolak \texttt{*.key.pem} sehingga kunci tidak masuk
ke layer, dan kunci diberikan saat eksekusi lewat berkas sementara yang di-mount
\textit{read-only}.

\paragraph{Kesalahan parameter.}
Setiap kunci divalidasi terhadap profilnya sebelum dipakai: RSA harus bermodulus tepat
3072 bit dengan $e = 65537$, ECDSA harus pada P-256 dengan titik SEC1 \textit{uncompressed},
dan Ed25519 mengikuti RFC 8410. Pada RSA-PSS, salt 32 byte dan MGF1-SHA-256 dipaksakan
saat verifikasi, sehingga tanda tangan PKCS\#1 v1.5 atau PSS dengan salt lain ditolak
(diuji pada \textit{interop} dengan OpenSSL). SHA-256 diterapkan satu kali oleh pustaka,
sehingga tidak terjadi \textit{double hashing}. Fungsi \texttt{verify} setiap profil tidak
pernah melempar galat: panjang atau nilai tidak sah menghasilkan \texttt{false}.

\paragraph{Kunci publik dari sumber tidak tepercaya.}
Keputusan kepercayaan hanya berasal dari \textit{trust store} yang disiapkan pengguna,
bukan dari bundle. Fingerprint kunci dihitung dari DER SPKI hasil ekspor ulang, sehingga
perbedaan \textit{line ending} atau pemenggalan PEM tidak mengubah \texttt{keyid}.

\subsection{Nonce dan Randomness (U14)}

Skenario U14 menandatangani manifest yang identik dua kali dengan kunci yang sama.
Hasilnya ditentukan oleh sumber randomness tiap algoritma (berkas
\texttt{test/nonce\_analysis\_test.exs}):

\begin{table}[H]
\centering
\begin{tabular}{@{}lll@{}}
\toprule
Profil & Randomness & Dua tanda tangan atas pesan sama \\
\midrule
ECDSA P-256 & nonce $k$ dari CSPRNG & berbeda \\
RSA-PSS & salt 32 byte dari CSPRNG & berbeda \\
Ed25519 & nonce dari hash kunci dan pesan (RFC 8032) & identik \\
\bottomrule
\end{tabular}
\caption{Perilaku tanda tangan untuk manifest yang identik.}
\end{table}

Pada ECDSA, 200 tanda tangan atas manifest yang sama menghasilkan 200 nilai $r$ yang
berbeda dan seluruhnya lolos verifikasi. Nilai $r$ yang berulang akan menandakan nonce
dipakai ulang. Pada RSA-PSS, 10 tanda tangan berbeda semuanya sah karena salt acak
mengubah blok yang ditandatangani. Pada Ed25519, nonce diturunkan secara deterministik
dari kunci dan pesan, sehingga tanda tangan identik dan tidak bergantung pada kualitas
generator acak saat menandatangani.

\paragraph{Dampak penggunaan ulang nonce ECDSA.}
Bila dua pesan $m_1 \neq m_2$ ditandatangani dengan $k$ yang sama, maka $r$ sama dan
\begin{align*}
s_i &= k^{-1}(z_i + r d) \bmod n, \quad z_i = \mathrm{SHA256}(m_i),\\
k &= (z_1 - z_2)(s_1 - s_2)^{-1} \bmod n,\\
d &= (s_1 k - z_1)\, r^{-1} \bmod n .
\end{align*}
Kunci privat $d$ pulih dari dua tanda tangan publik. Peragaan ini dijalankan pada kunci
uji dalam berkas yang sama dan memulihkan $k$ serta $d$ dengan tepat. Mitigasi pada
program ini adalah nonce dari CSPRNG OTP untuk setiap tanda tangan, tanpa seed tetap
(seed tetap hanya untuk \textit{test vector}). Alternatifnya adalah nonce deterministik
RFC 6979, yang diperbolehkan spesifikasi tetapi tidak dipakai karena pustaka OTP tidak
menyediakannya.

\subsection{\textit{Supply Chain}}

\begin{itemize}
  \item Seluruh \textit{action} GitHub dikunci per SHA commit, dan citra dasar Docker dikunci per
        digest.
  \item Dependensi program hanya pustaka OTP (\texttt{crypto}, \texttt{public\_key}),
        tanpa pustaka kripto eksternal. Dependensi pengembangan (ExDoc, Credo, Dialyxir)
        tidak masuk ke \textit{release}. CI memakai \texttt{mix deps.get --check-locked}
        dan \texttt{mix hex.audit} terhadap \texttt{mix.lock}.
  \item Izin workflow dibatasi \texttt{contents: read}. \textit{Checkout} memakai
        \texttt{persist-credentials: false}. CI berjalan pada \texttt{pull\_request}
        (bukan \texttt{pull\_request\_target}), sehingga PR dari \textit{fork} tidak
        menerima \textit{secret}. Run lama pada PR dibatalkan lewat \texttt{concurrency}.
  \item Token Docker Hub hanya ada pada \textit{environment} \texttt{nadina}, dipakai satu
        langkah \texttt{docker/login-action} pada job rilis setelah CI lulus.
  \item Image dipindai Trivy dan rilis gagal bila ada temuan HIGH atau CRITICAL.
        Pengecualian hanya lewat \texttt{.trivyignore} dengan alasan dan tanggal kedaluwarsa.
  \item Image berjalan sebagai pengguna non-root tanpa distribusi Erlang, dan tetap
        berfungsi dengan UID acak serta \texttt{--network none}.
\end{itemize}

\subsection{Rollback Attack}
Rancangan saat ini tidak mencegah \textit{rollback attack}. Bundle lama yang sah (tanda
tangan benar, kunci tepercaya, metadata cocok dengan masukan pengguna) tetap lolos bila
pengguna menyebut versi lama itu, karena tidak ada informasi tentang versi terbaru yang
pernah diterima. Pendeteksian membutuhkan \textit{trusted state} tambahan, misalnya
versi terendah yang diterima per paket yang disimpan di trust store, atau metadata
bertimestamp dan berkedaluwarsa seperti pada TUF.

\subsection{Keterbatasan}
% TODO: dilengkapi setelah seluruh fitur selesai (kompromi kunci dicabut lewat trust remove,
% kunci demo, tidak ada transparency log).
